 Si en la \autoref{def:grafm1} cambiamos  $A \subseteq \mathbb{R}$ por $A \subseteq \mathbb{R}^{n}$  y  $B \subseteq \mathbb{R}$ por $B \subseteq \mathbb{R}^{m}$ obtenemos de manera an\'aloga la siguiente definici\'on.    


\begin{definition}\textbf{Gráfica de un campo vectorial}
\label{def:grafica_campo_vec}
Sea $\campoVec{F} \colon A \subseteq \mathbb{R}^n \to \mathbb{R}^m$ un campo vectorial,  se define la \textit{gráfica de $\campoVec{F}$}, y se denota como $\operatorname{Graf}(\mathbf{f})$, al conjunto:
\[ 
\operatorname{Graf}(\campoVec{F}) = \{ (\mathbf{x}, \mathbf{y}) \in \mathbb{R}^{n} \times \mathbb{R}^{m} \mid \mathbf{x} \in A \land \mathbf{y} = \campoVec{F}(\mathbf{x}) \} 
\]
\end{definition}

\begin{obs} Siendo $\campoVec{F} \colon A \subseteq \mathbb{R}^n \to \mathbb{R}^m$ un campo vectorial entonces  desde la Definicón \ref{def:grafica_campo_vec} se desprende de manera inmediata que 
 \[
\operatorname{Graf}(\campoVec{F})   \subseteq  \mathbb{R}^{n+m}.
 \]
\end{obs}
Para el caso particular de un campo escalar obtenemos la siguiente definición
\begin{definition}\textbf{Gráfica de un campo escalar.} 
 \label{def:grafica}
 Sea \( f: A  \subseteq  \mathbb{R}^n \to  \mathbb{R}\),  se define la \textit{gráfica de $f$}, y se nota $\text{Graf}(f)$,  a
 \[
\operatorname{Graf}(f)=\{(\mathbf{x}, y) \in  \mathbb{R}^{n} \times \mathbb{R}  \;|\; \mathbf{x} \in A \land  y = f(\mathbf{x}) \}
 \]
\end{definition}


\begin{obs}  Siendo \( f: A  \subseteq  \mathbb{R}^n \to  \mathbb{R}\) un campo escalar   entonces  desde la Definicón \ref{def:grafica}  se desprende que  $$\text{Graf}(f) \subseteq  \mathbb{R}^{n+1}.$$  Con el objetivo de poder dibujar,  nos detendremos especialmente en el caso de gr\'aficas de  campos  $f: A\subseteq \mathbb{R}^2 \rightarrow \mathbb{R},$ cuya definici\'on es,  
 \[
\operatorname{Graf}(f)=\{(x,y,z)\in \mathbb{R}^3 \;|\;(x,y)\in A \land z=f(x,y) \}.
 \]
\end{obs}

 



